\documentclass[aspectratio=169,9pt]{beamer}

\usetheme{default}
\usefonttheme{professionalfonts}
\setbeamertemplate{navigation symbols}{}
\setbeamertemplate{footline}{\hfill\scriptsize\insertframenumber/\inserttotalframenumber\hspace{2mm}\vspace{1mm}}

\usepackage{amsmath,amssymb,bm,mathtools}
\usepackage{xcolor}
\usepackage[most]{tcolorbox}

\definecolor{deepblue}{RGB}{0,70,150}
\definecolor{softblue}{RGB}{232,242,255}
\definecolor{softgreen}{RGB}{235,250,238}
\definecolor{softorange}{RGB}{255,245,230}
\definecolor{softpurple}{RGB}{245,238,255}
\definecolor{darkred}{RGB}{160,35,25}
\definecolor{darkgreen}{RGB}{20,120,60}
\definecolor{darkpurple}{RGB}{95,45,150}

\setbeamerfont{frametitle}{series=\bfseries,size=\Large}
\setbeamercolor{frametitle}{fg=black}
\setbeamertemplate{itemize item}{\color{deepblue}$\blacktriangleright$}
\setlength{\abovedisplayskip}{3pt}
\setlength{\belowdisplayskip}{3pt}
\setlength{\jot}{1pt}

\newtcolorbox{bluebox}[1][]{colback=softblue,colframe=deepblue,boxrule=0.8pt,arc=1.5mm,left=1.5mm,right=1.5mm,top=0.8mm,bottom=0.8mm,#1}
\newtcolorbox{greenbox}[1][]{colback=softgreen,colframe=darkgreen,boxrule=0.8pt,arc=1.5mm,left=1.5mm,right=1.5mm,top=0.8mm,bottom=0.8mm,#1}
\newtcolorbox{orangebox}[1][]{colback=softorange,colframe=darkred,boxrule=0.8pt,arc=1.5mm,left=1.5mm,right=1.5mm,top=0.8mm,bottom=0.8mm,#1}
\newtcolorbox{purplebox}[1][]{colback=softpurple,colframe=darkpurple,boxrule=0.8pt,arc=1.5mm,left=1.5mm,right=1.5mm,top=0.8mm,bottom=0.8mm,#1}

\newcommand{\Evec}{\bm{E}}
\newcommand{\Dvec}{\bm{D}}
\newcommand{\nvec}{\bm{n}}
\newcommand{\xvec}{\bm{x}}
\newcommand{\epsSi}{\varepsilon_{\mathrm{Si}}}
\newcommand{\rhoDef}{\rho_{\mathrm{def}}}
\newcommand{\OmegaD}{\Omega}

\begin{document}

\begin{frame}{Module 2: Electrostatic map from radiation-induced charge}
\footnotesize
\vspace{-4mm}
\begin{bluebox}
\centering\small
\textbf{Goal:} convert dopants, mobile carriers, and charged radiation-induced defects into \(\phi(x,y,t)\) and \(\Evec(x,y,t)\).
\end{bluebox}

\vspace{1mm}
\begin{columns}[T,totalwidth=\textwidth]
\begin{column}{0.49\textwidth}
\textbf{Electrostatic core}
\begin{align*}
\boxed{\nabla\cdot\Dvec=\rho},\qquad
\boxed{\Dvec=\varepsilon\Evec},\qquad
\boxed{\Evec=-\nabla\phi}
\end{align*}
\begin{equation*}
\boxed{\nabla\cdot\left(\varepsilon\nabla\phi\right)=-\rho}
\end{equation*}
\small
This is the Poisson equation solved by Module 2.  Radiation does not directly appear as an electric field source; it appears through the charge density \(\rho\).
\end{column}

\begin{column}{0.49\textwidth}
\textbf{Space charge in irradiated silicon}
\begin{equation*}
\boxed{\rho=q\left(p-n+N_D^+-N_A^-\right)+\rhoDef}
\end{equation*}
\begin{equation*}
\boxed{\rhoDef=q\sum_{j=1}^{M} z_j C_j}
\end{equation*}
\small
\begin{itemize}\setlength\itemsep{0.2mm}
\item \(p,n\): mobile hole and electron densities.
\item \(N_D^+,N_A^-\): ionized donors and acceptors.
\item \(C_j,z_j\): concentration and charge number of defect species \(j\).
\end{itemize}
\end{column}
\end{columns}

\vspace{1mm}
\begin{orangebox}
\centering\scriptsize
\textbf{Module handoff:} \(C_j(x,y,t)\) from Module 3 changes \(\rho(x,y,t)\), which bends \(\phi\) and \(\Evec\), which then steers carrier transport in Module 5.
\end{orangebox}
\end{frame}

\begin{frame}{Module 2: Defect charge perturbs depletion and junction fields}
\footnotesize
\vspace{-4mm}
\begin{bluebox}
\centering\small
A charged defect map behaves like an additional space-charge distribution inside the device, shifting depletion structure and local junction fields.
\end{bluebox}

\vspace{1mm}
\begin{columns}[T,totalwidth=\textwidth]
\begin{column}{0.50\textwidth}
\textbf{Baseline plus radiation perturbation}
\begin{align*}
\phi &= \phi_0+\delta\phi,\qquad \rho=\rho_0+\delta\rho,\\[-0.5mm]
\delta\rho &= q\left(\delta p-\delta n+\sum_j z_j\delta C_j\right).
\end{align*}
For constant \(\varepsilon=\epsSi\), the perturbation obeys
\begin{equation*}
\boxed{\nabla^2\delta\phi=-\frac{\delta\rho}{\epsSi}}
\end{equation*}
\begin{equation*}
\boxed{\delta\Evec=-\nabla\delta\phi}
\end{equation*}
\small
Positive defect charge pulls field lines inward; negative defect charge reverses the local bending direction.
\end{column}

\begin{column}{0.48\textwidth}
\textbf{2D project equation}
\begin{equation*}
\boxed{
\frac{\partial^2\phi}{\partial x^2}
+\frac{\partial^2\phi}{\partial y^2}
=-\frac{q}{\epsSi}
\left[p-n+N_D^+-N_A^-+\sum_j z_j C_j\right]}
\end{equation*}
\small
\begin{itemize}\setlength\itemsep{0.2mm}
\item Charged traps alter the local electrostatic landscape.
\item The depletion boundary shifts because the local net space charge changes.
\item Neutral defects do not enter \(\rho\) directly, but still affect Module 5 through mobility loss and recombination.
\end{itemize}
\end{column}
\end{columns}

\vspace{1mm}
\begin{purplebox}
\scriptsize
\textbf{Physical interpretation:} Module 2 is the bridge from material damage to field distortion: defect charge modifies the internal electric field even before solving carrier transport.
\end{purplebox}
\end{frame}

\begin{frame}{Module 2: FEM weak form, boundary conditions, and outputs}
\footnotesize
\vspace{-4mm}
\begin{bluebox}
\centering\small
The numerical task is a 2D FEM Poisson solve with material regions, contacts, oxide interfaces, and defect-dependent charge density.
\end{bluebox}

\vspace{1mm}
\begin{columns}[T,totalwidth=\textwidth]
\begin{column}{0.50\textwidth}
\textbf{Weak form for finite elements}
\begin{equation*}
\boxed{
\int_{\OmegaD}\varepsilon\,\nabla\phi\cdot\nabla v\,d\Omega
=
\int_{\OmegaD}\rho\,v\,d\Omega
+\int_{\Gamma_N}\sigma_s v\,d\Gamma}
\end{equation*}
\small
for all test functions \(v\). Here \(\Gamma_N\) is a Neumann boundary and \(\sigma_s\) is a prescribed surface charge density.

\vspace{1mm}
\textbf{Typical boundary conditions}
\begin{align*}
\phi &= V_A \quad &&\text{on contact A},\\
\phi &= V_B \quad &&\text{on contact B},\\
-\nvec\cdot\varepsilon\nabla\phi &= \sigma_s \quad &&\text{on charged interfaces}.
\end{align*}
\end{column}

\begin{column}{0.48\textwidth}
\textbf{Outputs used by the rest of the project}
\begin{equation*}
\boxed{\phi(x,y,t)},\qquad
\boxed{\Evec(x,y,t)=-\nabla\phi(x,y,t)}
\end{equation*}
\small
\begin{itemize}\setlength\itemsep{0.3mm}
\item \textbf{Module 3:} defect charge state and migration can depend on electrostatic potential.
\item \textbf{Module 5:} carrier drift currents use \(\Evec\):
\end{itemize}
\begin{align*}
\bm{J}_n&=q\mu_n n\Evec+qD_n\nabla n,\\[-0.5mm]
\bm{J}_p&=q\mu_p p\Evec-qD_p\nabla p.
\end{align*}
\begin{itemize}\setlength\itemsep{0.3mm}
\item \textbf{Module 6:} iterates \(C_j,T,n,p,\phi\) until maps are self-consistent.
\end{itemize}
\end{column}
\end{columns}

\vspace{1mm}
\begin{greenbox}
\centering\scriptsize
\textbf{Validation check:} changing the sign or magnitude of \(z_jC_j\) should produce the corresponding local bending of \(\phi\), \(\Evec\), depletion width, and carrier current paths.
\end{greenbox}
\end{frame}

\end{document}
